% !PRINTED 276
\weberpage{276}
Diese Gruppe ist reell, wenn $n$ ein Theiler von $1/2\,(p-1)$ ist, weil
dann und nur dann die Exponenten von $\gamma$ durch $p+1$ theilbar
sind. Die Gruppe ist in $\Gamma_p$ und zugleich in $\Lambda_p$ (\S.\,69)
enthalten, wenn $\gamma^{\frac{p^2-1}{2n}}$, $\gamma^{-\frac{p^2-1}{2n}}$
conjugirt sind, wenn also
\[
\gamma^{-\frac{p^2-1}{2n}} = \gamma^{p\frac{p^2-1}{2n}}
\quad \text{oder} \quad
\gamma^{(p+1)\frac{p^2-1}{2n}} = 1
\]
ist, d. h. wenn $n$ ein Theiler von $1/2\,(p+1)$ ist. Dies stimmt mit
\S.\,68 überein, wonach andere cyklische Gruppen, als solche, deren
Grad gleich $p$ oder ein Theiler von $1/2\,(p-1)$ oder $1/2\,(p+1)$
ist, nicht vorkommen. In den übrigen Polyëdergruppen ist $h = 3$, und
wir haben:

II. \gesperrt{Diëdergruppen} vom Grade $2m$,\quad
$\begin{aligned} \nu &= \nu' = 2,\ \nu'' = m \\ \mu &= \mu' = m,\ \mu'' = 2 \end{aligned}$
\[
D_m = \begin{pmatrix} \gamma^{\frac{p^2-1}{2m}\lambda}, & 0 \\[2pt] 0, & \gamma^{-\frac{p^2-1}{2m}\lambda} \end{pmatrix},
\begin{pmatrix} 0, & \gamma^{\frac{p^2-1}{2m}\lambda} \\[2pt] -\gamma^{-\frac{p^2-1}{2m}\lambda}, & 0 \end{pmatrix},
\quad \lambda = 0, 1, \ldots, m-1.
\]
Diese Gruppe ist in $L_p$ enthalten, wenn $p \equiv 1 \pmod{2m}$, und
in $\Gamma_p$, wenn $p \equiv -1 \pmod{2m}$, wie im Falle I.

Hier ist die in \S.\,55 gegebene zweite Darstellung der Diëdergruppe
gewählt, in der die Unterscheidung der Fälle etwas einfacher wird, als
bei der ersten.

III. \gesperrt{Tetraëdergruppe,}\quad
$\begin{aligned} \nu &= 2,\ \nu' = 3,\ \nu'' = 3,\ n = 12 \\ \mu &= 6,\ \mu' = 4,\ \mu'' = 4. \end{aligned}$

Um diese Gruppe darzustellen, bemerken wir, dass im Körper
$\mathfrak{E}_{2,p}$ immer eine $8^{\mathrm{te}}$ Einheitswurzel
$\gamma^{\frac{p^2-1}{8}}$ existirt. Wir erhalten also nach \S.\,56 die
drei erzeugenden Substitutionen
\[
\Theta = \begin{pmatrix} \gamma^{\frac{p^2-1}{4}}, & 0 \\[2pt] 0, & \gamma^{-\frac{p^2-1}{4}} \end{pmatrix},
\quad
\psi = \begin{pmatrix} 0, & \gamma^{\frac{p^2-1}{4}} \\[2pt] -\gamma^{-\frac{p^2-1}{4}}, & 0 \end{pmatrix},
\quad
\psi_1 = \psi\,\Theta = \begin{pmatrix} 0, & 1 \\ -1, & 0 \end{pmatrix}
\]
\[
\chi = \begin{pmatrix} \dfrac{1}{\sqrt{-2}}\,\gamma^{\frac{p^2-1}{8}}, & \dfrac{1}{\sqrt{-2}}\,\gamma^{\frac{p^2-1}{8}} \\[8pt] \dfrac{1}{\sqrt{-2}}\,\gamma^{-\frac{p^2-1}{8}}, & -\dfrac{1}{\sqrt{-2}}\,\gamma^{-\frac{p^2-1}{8}} \end{pmatrix}.
\]
1) Ist $p \equiv 1 \pmod 4$, so ist $\gamma^{\frac{p^2-1}{4}}$ reell,
und $\gamma^{\frac{p^2-1}{8}}$ ist reell oder rein imaginär, je
nachdem $p \equiv 1$ oder $\equiv 5 \pmod 8$ ist.
